\section{Provas}

\subsection{Prova 1}

\begin{enumerate}[left=0.5cm, align=left, nosep]
    \item Seja $p$ primo e $n$ um número inteiro.
    \begin{enumerate}[label=\alph*), left=0.5cm, align=left, nosep]
        \item Mostre por indução que $p$ é múltiplo de $n^p - n$ para $n \geq 0$.  
        \item Mostrar que o resultado do item anterior vale para $n \in \mathbb{Z}$ (incluindo negativos).
    \end{enumerate}

    \item Dado os seguintes polinômios $X^7 - 3X^5 + 2X^4$ e $X^5 + X^4 - 2X^3 - X^2 - X + 2$:
    \begin{enumerate}[label=\alph*), left=0.5cm, align=left, nosep]
        \item Encontre o m.d.c entre os polinômios.
        \item Escrever o polinômio na forma fatorada pelas raízes.
        \item Encontrar uma combinação linear que é igual à ...
    \end{enumerate}

    \item Mostre que: 
    \begin{enumerate}[label=\alph*), left=0.5cm, align=left, nosep]
        \item Existem infinitos números primos usando o argumento de Euclides. 
        \item Existem infinitos números primos gaussianos.
    \end{enumerate}

    \item Dada a seguinte operação...
    \begin{enumerate}[label=\alph*), left=0.5cm, align=left, nosep]
        \item A operação é associativa? Justifique. 
        \item A operação é comutativa? Justifique.
    \end{enumerate}
\end{enumerate}

\subsection{Prova 2}

\begin{enumerate}[left=0.5cm, align=left, nosep]
    \item Sejam $G$ e $G'$ grupos e seja $f: G \to G'$ uma função bijetiva tal que $f(xy) = f(x)f(y)$ para todos $x, y \in G$.
    \begin{enumerate}[label=\alph*), left=0.5cm, align=left, nosep]
        \item Seja o conjunto $S = \{ x \in G : f(x) = e' \}$, onde $e'$ é a identidade de $G'$. Verifique se $S$ é um subgrupo de $G$.
        \item Seja o conjunto $H = \{ x \in G : f(x) = x \}$. Verifique se $H$ é um subgrupo de $G$ ou $G'$.
    \end{enumerate}

    \item Considere o grupo multiplicativo $(\mathbb{Z}/m\mathbb{Z})^{\times}$ do anel $(\mathbb{Z}/m\mathbb{Z})$.
    
    \begin{enumerate}[label=\alph*), left=0.5cm, align=left, nosep]
        \item Os grupos multiplicativos $(\mathbb{Z}/32\mathbb{Z})^\times$ e $(\mathbb{Z}/34\mathbb{Z})^\times$ são isomorfos? Justifique.
        \item Os anéis $\mathbb{Z}/32\mathbb{Z}$ e $\mathbb{Z}/34\mathbb{Z}$ são isomorfos? Justifique.
    \end{enumerate}

    \item Mostre que:
    \begin{enumerate}[label=\alph*), left=0.5cm, align=left, nosep]
        \item Se $m \geq 1$ não é primo, então o anel $\mathbb{Z}/m\mathbb{Z}$ não é um corpo.
        \item Se $p$ é primo, então o anel $\mathbb{Z}/p\mathbb{Z}$ é um corpo.
    \end{enumerate}

    \item Sejam $X = \{1, 2, 3, 4, 5\}$ e o grupo G das bijeções $X \rightarrow X$ sob a composição $\circ$ de funções. Considere 
    os elementos $\sigma, \tau \in G$ definidos por: 
    \[ \sigma(1) = 2,\ \sigma(2) = 3,\ \sigma(3) = 4,\ \sigma(4) = 5,\ \sigma(5) = 1 \] 
    \[ \tau(1) = 1,\ \tau(2) = 5,\ \tau(3) = 4,\ \tau(4) = 3,\ \tau(5) = 2 \] 
    \begin{enumerate}[label=\alph*), left=0.5cm, align=left, nosep] 
        \item Determine explicitamente o elemento inverso de $\tau \circ \sigma$ em $G$. 
        \item É verdade que $\tau \circ \sigma = \sigma^4 \circ \tau$? Justifique.
    \end{enumerate}
\end{enumerate}

\subsection{Prova 3}

\begin{enumerate}[left=0.5cm, align=left, nosep]
    \item Seja $A$ um domínio de integridade e $S$ o conjunto dos elementos não nulos de $A$.Demonstre que a relação sobre 
    $A \times S$ definida por $(a,s)$ está relacionado a $(a',s')$ se e somente se, existe um $s'' \in S$ tal que 
    $s''(s'a - sa') = 0$ é uma relação de equivalência

   \item Sobre $\mathbb{C}$ considere as relações de equivalência $R$ e $S$ definidas por: $(a,b)$ pertence a $R$ ou $S$ 
   de acordo quando $a - b$ pertence a $\mathbb{Z}$ ou $\mathbb{Z}[i]$.Seja, $f: \mathbb{C} \rightarrow \mathbb{C}$ e 
   $f_* : \mathbb{C}/R \rightarrow \mathbb{C}/S$ definidas por $f(x) = x$ e $f_*((x \mod R)) = (f(x) \mod S)$. 
    \begin{enumerate}[label=\alph*), left=0.5cm, align=left, nosep] 
        \item Verifique que a definição de $f_*$ não depende da escolha de representantes. 
        \item A função $f_*$ é um isomorfismo entre grupos aditivos com respeito às estruturas quocientes usuais ? Justifique. 
    \end{enumerate}

    \item
    \begin{enumerate}[label=\alph*), left=0.5cm, align=left, nosep] 
        \item Verifique que a aplicação $f: \mathbb{Z}[X] \rightarrow \mathbb{Z}$ entre anéis unitários dada por: 
        $f(a_nX^n + \ldots + a_1X + a_0 ) = a_0$ é um homomorfismo. 
        \item Descreva o núcleo $K$ de $f$ e a imagem de $f$. 
        \item O anel quociente $\mathbb{Z}[X]/K$ é um domínio de integridade ? é um corpo ? Justifique. 
    \end{enumerate}

    \item Considere o grupo $G_n$ das simetrias de um polígono regular de $n$ lados com a operação de composição e o 
    grupo multiplicativo $(\mathbb{Z}/m\mathbb{Z})^{\times}$ do anel $(\mathbb{Z}/m\mathbb{Z})$. Justifique quais dos 
    seguintes grupos são isomorfos a $\mathbb{Z}/6\mathbb{Z}$.
    \begin{enumerate}[label=(\roman*), left=0.5cm, align=left, nosep]
        \item $\mathbb{Z}/2\mathbb{Z} \times \mathbb{Z}/3\mathbb{Z}$
        \item $(\mathbb{Z}/7\mathbb{Z})^{\times}$
        \item $(\mathbb{Z}/14\mathbb{Z})^{\times}$
        \item $G_3$
        \item $G_6$
    \end{enumerate}
    Justifique quais deles são isomorfos a $\mathbb{Z}/6\mathbb{Z}$.
\end{enumerate}

\subsection{Prova Substitutiva}

\begin{enumerate}[left=0.5cm, align=left, nosep]
    \item Mostre por indução que, para todo $n \geq 1$, $n^3 - 3n^2 - n$ é múltiplo de 6.

    \item Considere o grupo $G_n$ das simetrias de um polígono regular de $n$ lados com a operação de composição e o 
    grupo multiplicativo $(\mathbb{Z}/m\mathbb{Z})^{\times}$ do anel $(\mathbb{Z}/m\mathbb{Z})$. Justifique quais dos 
    seguintes grupos são isomorfos a $\mathbb{Z}/6\mathbb{Z}$.
    \begin{enumerate}[label=(\roman*), left=0.5cm, align=left, nosep]
        \item $\mathbb{Z}/2\mathbb{Z} \times \mathbb{Z}/3\mathbb{Z}$
        \item $(\mathbb{Z}/7\mathbb{Z})^{\times}$
        \item $(\mathbb{Z}/14\mathbb{Z})^{\times}$
        \item $G_3$
        \item $G_6$
    \end{enumerate}

    \item
    \begin{enumerate}[label=(\roman*), left=0.5cm, align=left, nosep]
        \item Mostre que todo domínio de integridade também é um corpo.
        \item Não me lembro ao certo. É verdade que todas as raízes do polinômio $X^2 - 5X$ estão em $\mathbb{Z}/6\mathbb{Z}$? Justifique.
    \end{enumerate}

    \item Quais das seguintes estruturas são corpos? Justifique. Considere $[p(x)]$ o ideal gerado pelo polinômio $p(x)$(Não lembro 
    questão ao certo).
    \begin{enumerate}[label=(\roman*), left=0.5cm, align=left, nosep]
        \item $\mathbb{Z}[X]/(X^2 - 1)$
        \item $\mathbb{R}[X]/(X^2 + 1)$
        \item $\mathbb{C}[X]/(X^2 + 3)$
        \item $\mathbb{Q}[X]/(X^2 + 2)$
        \item $\mathbb{R}[X]/(\text{?})$
    \end{enumerate}

    \item Verifique quais das seguintes operações são relações de equivalência:
    \begin{enumerate}[label=(\roman*), left=0.5cm, align=left, nosep]
        \item $\equiv (mod \ m)$
        \item $\leq$
    \end{enumerate}

\end{enumerate}
